% Transcription — fonds Alexandre Grothendieck, shelfmark 130, batch 2 (pages 21-40).
% Established from the facsimile archives/batches/130/batch-02.pdf (a mirror of
% https://grothendieck.umontpellier.fr/130.pdf), per the skill
% transcribe-grothendieck. The batch file carries no cover sheet: PDF page N
% is page 20+N of the folder (the Montpellier footer prints « 21 » ... « 40 »),
% and the archivists' pencil numbers agree wherever legible (« 21 », « 23 »,
% « 25 », « 27 », « 29 », « 31 », « 33 » ... « 40 »).
%
% Pass: Opus 5.5 (claude-opus-5-5), 2026-09-23 — first pass, unchecked against the pages by a human.
%
% Thirteen of the twenty pages are transcribed. Absent: pages 22, 24, 26, 28,
% 30, 32, typed pages of an unrelated duplicated typescript with nothing in his
% hand; page 33, a blank leaf. He did not paginate these leaves. The \section
% headings are bracketed where they are ours; the one before page 34,
% « Quasi-fonctions de Néron-Weil », and the \subsection headings are his own
% underlined titles. The folder title is the archivists', unbracketed.
%
% What the batch is. Two runs, undated.
%   21-31  Odd pages only, landscape half-sheets: the local pairing
%          <C_η, D_η>_v of cycles algebraically equivalent to 0 on the
%          generic fibre of X -> S (S a DVR spectrum), compared with the
%          intersection number i(C.D) on a regular model — the notation
%          <.,.>_v and i(C.D) is set up in batch 1 (page 16). Page 21 states
%          the global form -<D_η,C_η> = Σ_v <D_η,C_η>_v = Σ_v (D.C)_v (S with a
%          set M of discrete valuations satisfying the product formula), then
%          the local hypotheses and the boxed <C_η,D_η>_v = i(C.D). Page 23 is
%          a scratch sheet on Néron models (boxed: a 2-dimensional local ring
%          of depth 2 with A/fA regular; products A_η = A_1η × A_2η,
%          closures, isogenies B_η -> A_η -> B_η with composite N). Page 25:
%          the case D = div f, <C_η, div f_η>_v = i(C.div f) = v(f_η(C_η)), and
%          the choice of a line bundle L with a rational section. Pages 27-29:
%          « a) » f smooth, via the Albanese map α : X -> A (Néron model),
%          D = α*(D') + div(f), reducing to the div(f) formula of page 25;
%          page 31: « b) Cas général », α : X' = X - x_0 -> A, broken off.
%   34-40  « Quasi-fonctions de Néron-Weil »: K with a « proper » set M_K of
%          absolute values, real functions φ on X(K̄) × M_K̄, Galois
%          invariance, M_K-bounded functions and subsets E(U; x_1..x_n; c),
%          admissible functions, the example φ_f(P,v) = v(f(P)) for f a unit,
%          reduced to G_m; « Faisceautisation »: admissible Γ-invariant
%          functions form a presheaf, claimed a sheaf; page 40 breaks off
%          mid-sentence (batch 3 continues).
%
% Notation fixed for the batch: the closure of K as \overline{K}; his hatted
% capital for the set of quasi-functions read as \widehat{\Phi} (the glyph on
% pages 35-36 is a looped upright stroke; page 36 draws it most like a Φ),
% and the same letter underlined as \underline{\Phi}_X for the presheaf of
% page 40; Néron models written with an underbar as \underline{A}. Words he
% underlines are set in \emph.
\documentclass[11pt,a4paper]{article}
% Le préambule commun à toutes les transcriptions du fonds Grothendieck.
%
% Il tient en une idée : **l'incertitude se compose, elle ne se résout pas.**
% Une transcription qui lisse un mot douteux a détruit la seule chose pour
% laquelle on la faisait — un lecteur doit pouvoir savoir, à chaque endroit,
% ce qui était sur le feuillet et ce qui a été ajouté. Les six macros qui
% suivent sont donc l'appareil critique, et non de la décoration.
%
% Les mêmes commandes sont reconnues par `scripts/render.mjs`, qui produit la
% vue de lecture du site. Une macro ajoutée ici doit l'être aussi là-bas,
% faute de quoi le rendu échouera bruyamment — ce qui est voulu : un rendu
% silencieusement faux est pire qu'un rendu absent.

\ProvidesPackage{grothendieck}[2026/08/08 Transcriptions du fonds Grothendieck]

\usepackage{amsmath,amssymb,amsthm}
\usepackage{mathrsfs}
\usepackage{stmaryrd}
% Les diagrammes commutatifs sont partout dans ce fonds : c'est la langue
% même de la géométrie algébrique de Grothendieck, et une transcription qui
% les rendrait en prose ne transcrirait rien.
\usepackage{tikz-cd}
% Français par défaut — c'est la langue des feuillets — mais l'anglais est
% chargé aussi : la traduction et le résumé sont des documents anglais, et la
% typographie française (espaces avant les deux-points) y serait une faute.
% Ces documents basculent par \selectlanguage{english} en tête de corps.
\usepackage[english,french]{babel}
\usepackage{fontspec}
\usepackage{geometry}
\geometry{a4paper,margin=25mm,marginparwidth=28mm}
\usepackage{marginnote}
\usepackage{xcolor}
% ulem, not soul. `soul`'s \st and \ul are built on a character-by-character
% scan that XeLaTeX's font machinery breaks: the document fails with an
% undefined control sequence rather than degrading. ulem does the same two jobs
% and is Unicode-engine safe. `normalem` stops it hijacking \emph, which the
% transcriptions use for Grothendieck's own emphasis.
\usepackage[normalem]{ulem}
\usepackage{hyperref}
\hypersetup{colorlinks=true,linkcolor=black,urlcolor=black,pdfborder={0 0 0}}

% --- Métadonnées du lot -----------------------------------------------------
% Reprises telles quelles par le script de rendu, qui les affiche en tête de
% la vue de lecture. Elles fixent ce que le fichier prétend couvrir : sans
% elles, une transcription flotte sans point d'attache dans un fonds de seize
% mille feuillets.

\newcommand{\folder}[1]{\gdef\grothFolder{#1}}
\newcommand{\batch}[1]{\gdef\grothBatch{#1}}
\newcommand{\leaves}[2]{\gdef\grothFirst{#1}\gdef\grothLast{#2}}
\newcommand{\foldertitle}[1]{\gdef\grothFoldertitle{#1}}
\gdef\grothFolder{?}\gdef\grothBatch{?}\gdef\grothFirst{?}\gdef\grothLast{?}\gdef\grothFoldertitle{}

% --- Le résumé --------------------------------------------------------------
%
% Il ouvre la lecture modernisée, sous le titre « Résumé », et il est composé
% en petit. Ce n'est pas un ornement : c'est le seul passage du document qui
% ne soit pas des mathématiques, et un lecteur doit pouvoir le reconnaître
% avant de l'avoir lu — puis le sauter s'il connaît déjà le sujet, ou s'y
% arrêter s'il ne le connaît pas. Le filet qui le suit marque la reprise du
% corps.

\newenvironment{resume}
  {\small\setlength{\parskip}{0.45em}}
  {\par\medskip\hrule\medskip}

% --- Mots-clés ---------------------------------------------------------------
%
% En anglais, en clôture du résumé de la lecture modernisée : le vocabulaire
% moderne sous lequel chercher ce que les feuillets construisent. Le manifeste
% du site (`npm run manifest`) les extrait pour étiqueter la cote entière —
% cette ligne est la source unique des tags, il n'y a pas de fichier à côté.
\newcommand{\keywords}[1]{\par\smallskip\noindent
  {\footnotesize\textsc{Keywords} --- \textit{#1}}\par}

% --- L'appareil critique ----------------------------------------------------

% Le numéro de feuillet, en marge. C'est l'ancre qui permet de régler un
% désaccord sur une formule en regardant *un* feuillet plutôt qu'une chemise
% de six cents. Le site s'en sert aussi pour tourner les pages du fac-similé
% quand on fait défiler la transcription.
\newcommand{\leaf}[1]{%
  \marginnote{\footnotesize\color{gray!70}\textsf{#1}}%
  \label{leaf:#1}\ignorespaces}

% Illisible : on ne devine pas. Un mot inventé qui se lit comme les autres est
% le pire résultat possible.
\newcommand{\ill}{\textcolor{red!60!black}{[\,\ldots\,]}}

% Lecture douteuse : le mot est donné, et signalé comme incertain.
\newcommand{\uncertain}[1]{\uline{#1}}

% Ajout de l'éditeur — une lettre manquante, un mot sous-entendu.
\newcommand{\add}[1]{\textcolor{blue!50!black}{[#1]}}

% Biffé par Grothendieck lui-même. Ce qu'il a rayé fait partie du document :
% une rature dit souvent où la pensée a changé de direction.
\newcommand{\struck}[1]{\sout{#1}}

% Note du transcripteur, sur l'état matériel du feuillet.
\newcommand{\note}[1]{\par\smallskip{\small\color{gray!60!black}\textit{[#1]}}\par\smallskip}

% Note marginale de Grothendieck, distincte du corps du texte.
\newcommand{\marginal}[1]{\par\smallskip\noindent\hspace{1em}%
  {\small\color{gray!60!black}#1}\par\smallskip}

% --- Titre ------------------------------------------------------------------

\AtBeginDocument{%
  \begin{center}
    {\large\bfseries Fonds Alexandre Grothendieck --- cote n\textsuperscript{o}~\grothFolder}\\[2pt]
    {\itshape\grothFoldertitle}\\[4pt]
    {\small feuillets \grothFirst--\grothLast\ (lot \grothBatch)}\\[2pt]
    {\footnotesize\color{gray!60!black}Transcription établie d'après le fac-similé numérisé
     par l'Université de Montpellier.\\
     Les lectures incertaines sont soulignées, les illisibles notées [\,\ldots\,],
     les ajouts entre crochets.}
  \end{center}
  \vspace{1em}}

\endinput


\folder{130}
\batch{2}
\pages{21}{40}
\dating{[vers 1971]}
\watermark{Édition de démonstration}
\foldertitle{Hauteurs et symboles locaux : notes manuscrites (s.d.).}

\begin{document}

\section{[Accouplements locaux $\langle C_\eta, D_\eta \rangle_v$ et nombre d'intersection $i(C.D)$]}
\note{titre de l'éditeur, entre crochets. Les notations
$\langle C_\eta, D_\eta \rangle_v$ et $i(C.D)$ sont posées au lot 1
(feuillet 16). Les feuillets de cette suite sont des demi-feuilles tenues
en largeur ; seuls les feuillets impairs portent sa main.}

\page{21}

\begin{tikzcd}
X \arrow[d, "f"] \\
S
\end{tikzcd}

\[
  -\langle D_\eta, C_\eta \rangle = D.C
\]
si $D_\eta$ et $C_\eta$ alg. équiv. à 0 ;
\[
  -\langle D_\eta, C_\eta \rangle
  = \sum_v \langle D_\eta, C_\eta \rangle_v = \sum_v (D.C)_v
\]
\note{le second membre est écrit sous le premier, relié par un signe $=$
vertical.}

$f$ projectif, à fibres géom. \uncertain{irréductibles}, fibre générique
lisse, $X$ régulier ; $S$ \uncertain{muni d'un ensemble} $M$ de
\uncertain{val.} discrètes satisfaisant la formule du produit.
\note{ces hypothèses (situation globale) sont soulignées et séparées du reste
du feuillet par un long trait.}

$S$ local complet, corps résiduel alg. clos, $X \xrightarrow{f} S$
projectif, $X_\eta$ lisse, $X_0$ lisse sauf \ill{} en un nombre fini de pts,
ces derniers étant \ill{} \ill{}, $X$ \emph{\uncertain{régulier}}
\uncertain{en ces pts} ;
\note{la flèche $X \to S$ est tracée de droite à gauche ; « en ces pts » est
écrit au-dessus de la ligne.}

$C$, $D$ sur $X$, \struck{se \ill{}} : pas de comp. verticale ;
\struck{$C.D$} $\operatorname{Supp} C \cap \operatorname{Supp} D \subset X_0$
\quad $=$ \uncertain{se coupent avec la bonne dimension} ;
\note{les deux conditions sont réunies par une accolade.}

$C_\eta$, $D_\eta$ alg. équiv. à 0 sur $X_\eta$ :
\[
  \boxed{\ \langle C_\eta, D_\eta \rangle_v = i(C.D)\ }
\]

\ill{} \ill{} \ill{}
\note{trois mots soulignés au bas du feuillet, illisibles.}

\page{23}
$A$ anneau local \uncertain{noeth.} \uncertain{réduit} de dim 2, de
profondeur $= 2$, $f \in \mathfrak{m}_A$ tel que $A/fA$ soit régulier.
\note{encadré en haut du feuillet ; « noeth. » est ajouté au-dessus de la
ligne.}

\note{Le reste du feuillet est une page de brouillon sur les modèles de
Néron, sans phrase suivie ; on en donne les éléments dans l'ordre de la page.
Le soulignement sous une lettre ($\underline{A}$, $\underline{B}$) est le
sien.}

$\underline{X}$ \note{un trait vertical descend sous la lettre, qui est
barrée d'un trait oblique.}

$\underline{A} \supset B$ ; \quad $B \times B'$ ; \quad
$A_\eta = A_{1\eta} \times A_{2\eta}$, \quad $\underline{A}$ ok.

$B_1 = \overline{A_{1\eta}}$, \quad $B_2 = \overline{A_{2\eta}}$ ; \quad
$\underline{A}$ ; \quad
$\underline{A} \to \overline{A_{1\eta}}$, \quad
$\underline{A} \to \overline{A_{2\eta}}$.
\note{un grand $X$ isolé au milieu du feuillet.}

\begin{tikzcd}
& \underline{T} \\
\underline{A}^{\circ} & \underline{B} \arrow[l]
\end{tikzcd}

\begin{tikzcd}
A_\eta & B_\eta \arrow[l] \\
& T_\eta \arrow[u]
\end{tikzcd}

\note{dans le premier schéma, $\underline{T}$ est écrit directement au-dessus
de $\underline{B}$, sans flèche ; l'exposant de $\underline{A}$ est un petit
point, lu $\circ$.}

\begin{tikzcd}
\underline{A} \arrow[r, bend left=20] & \underline{B_1 \times B_2} \arrow[l, bend left=20] \\
B_1 \times B_2 \arrow[u]
\end{tikzcd}

\note{les deux flèches entre $\underline{A}$ et $\underline{B_1 \times B_2}$
sont superposées, en sens contraires ; le $B_1$ est surchargé. Au-dessous, un
croquis : deux traits verticaux, l'un simple, l'autre hachuré, posés sur un
trait horizontal.}

\begin{tikzcd}
B_\eta & A_\eta \arrow[l, "\beta"'] & B_\eta \arrow[l, "\alpha"'] \arrow[ll, bend right=40, "N"'] \\
\underline{B} & \underline{A} \arrow[l] \arrow[u] & \underline{B} \arrow[l]
\end{tikzcd}

$N_\eta$
\note{écrit sous ce diagramme.}

\begin{tikzcd}
0 \arrow[d, leftrightarrow] \\
\underline{A} \\
\underline{B} \arrow[u] \\
C \arrow[u]
\end{tikzcd}

\begin{tikzcd}
0 \\
A \arrow[d] \\
B \arrow[d] \\
C
\end{tikzcd}

\note{deux colonnes côte à côte sur le feuillet, la seconde d'un trait plus léger ; la
première flèche de la première colonne porte une pointe aux deux bouts,
lecture douteuse ; en tête de la seconde, un « $0$ » ou un « $i$ », lecture douteuse.}

$A \subset B$ ; \quad $\overline{A}$ ; \quad $B \to C$ ; \quad
$\underline{A} \to S$.

\page{25}
\[
  \langle C_\eta, \operatorname{div} f_\eta \rangle_v
  = i(C.\operatorname{div} f)
  = v\bigl(f_\eta(C_\eta)\bigr)
\]
\note{dans $i(C.\operatorname{div} f)$, un $f_\eta$ est biffé et remplacé
par $f$ au-dessus de la ligne ; la suite de la ligne est biffée.}

\begin{tikzcd}
X_\eta & X \arrow[d, no head] & X' \arrow[r, hook] \arrow[d, no head] & J \\
\eta & S & S
\end{tikzcd}

$S$ \uncertain{spectre d'anneau de val. discrète}.

\begin{tikzcd}
X' \arrow[r, hook] & \widetilde{X''} \arrow[r, "\mathrm{birat.}"] & X'' \arrow[r, hook] & J
\end{tikzcd}

$X''_0 = X_0$
\note{au-dessus de la flèche $\widetilde{X''} \to X''$ : « \ill{}
\emph{birat.} \ill{} \ill{} » ; une flèche en pointillé part de $X'$ vers
la droite, avec un mot illisible. La flèche entre $X''$ et $J$ est tracée
de droite à gauche, crochet du côté de $X''$. Sous la ligne, des arcs
esquissent les fibres.}

$x \in X_\eta(k(\eta))$ \quad $\exists\, L$ faisceau loc. libre sur $X$ et
une section rationnelle \uncertain{dont le} diviseur sur $X_\eta$
\uncertain{est $x$} : $(*)$

En normalisant, on peut supposer que cette section est définie partout, \ill{}
en \struck{$x \in X$} le pt \uncertain{critique} $x_0$, et \uncertain{que}
\struck{\ill{}} son diviseur n'a pas de composante verticale. Ceci
\ill{} \uncertain{pour} la section donnée \ill{} un diviseur.
\note{le mot après « partout » est souligné. Au bas du feuillet, deux petits
croquis : un cercle traversé d'un arc épais et d'un axe vertical, un cercle
centré sur un axe vertical.}

\page{27}
\uncertain{a)} Cas où $f$ est lisse, alors $X_\eta/K$ a un Albanese
$A_\eta/K$ et Picard°° $A^*_\eta/K$, \struck{\ill{}} \ill{} \ill{} $X$ sur
$S$, \uncertain{cf. th. de Weil}, \ill{}
\note{le repère en tête (« a) ») est peu formé ; le corps de base est lu $K$.
Les deux lignes sous « Albanese » sont en partie biffées (on y lit $A$,
$A^*$ et « Albanese ») ; une insertion soulignée au-dessus de la ligne est
illisible.}

\begin{tikzcd}
X \arrow[rr, "\alpha"] \arrow[dr, no head] & & A \arrow[dl, no head] \\
& S &
\end{tikzcd}

\struck{\ill{} de $X$ sur $A_\eta$ \ill{}}

$D'_\eta$ tel que $\alpha^*(D'_\eta) \sim_\ell D_\eta$, d'où
\[
  \alpha^*(D'_\eta) = D_\eta + \operatorname{div}_{X_\eta}(f)
\]
d'où un \ill{} prenant les adhérences
\note{ce bloc est encadré et annulé par deux traits obliques.}

\struck{$D = \overline{D_\eta} = \alpha^*(\overline{D'_\eta})$}

Supposons d'abord que \ill{}
\[
  D = \alpha^*(D')
\]
$D'$ un diviseur sur $A$ \uncertain{sans} \struck{\ill{}} tel que
$\alpha^*(D')$ défini. Posons
\[
  C' = \alpha_*(C) .
\]
Alors on a
\[
  C'.D' = \alpha_*(C).D' = \alpha_*\bigl(C.\alpha^*(D')\bigr)
  = \alpha_*(C.D)
\]
d'où
\[
  i(C.D) = i\bigl(\alpha_*(C.D)\bigr) = i(C'.D')
\]

\page{29}
et notre formule dans ce cas se réduit à :
\[
  \boxed{\ i\bigl(\overline{C'_\eta}.\overline{D'_\eta}\bigr)
  = \langle C'_\eta, D'_\eta \rangle_v\ }
\]
\struck{si $C'$, $D'$ sans \ill{}}

\struck{\ill{}} \uncertain{Dans le} cas général, on \uncertain{voit}
aussitôt \uncertain{qu'on} \ill{}
\[
  D_\eta \sim \alpha^*(D'_\eta) ,
\]
où $D'_\eta$ est un diviseur sur $A_\eta$ tel que
\[
  D'_\eta \not\supset \alpha_\eta(X_\eta) \quad\text{et}\quad
  \overline{D'_\eta} \not\supset \alpha_0(X_0)
\]
\note{sur le premier $D'_\eta$, une barre d'adhérence est biffée ; entre
$\overline{D'_\eta}$ et le second $\not\supset$, un petit rond.}

[\uncertain{Lemme} \ill{} : \ill{}]

d'où
\[
  D = \alpha^*(D') + \operatorname{div}(f) ,
\]
où $D'$ est tel que $\alpha^*(D')$ soit défini et sans composante
verticale, et $f \in R^*(X)$. \struck{Donc}

On est donc ramené à prouver la formule
\[
  \boxed{\ \langle C_\eta, \operatorname{div} f_\eta \rangle_v
  = i(C.\operatorname{div} f)\ }
  = v\bigl(f_\eta(C_\eta)\bigr)
\]
valable si $\operatorname{div}(f)$ est sans composante horizontale.
\note{le dernier membre est écrit sous le premier, relié par un signe $=$
vertical ; la condition est écrite en marge droite de la formule.}

\page{31}
b) Cas général. Soit $A$ le modèle de Néron de $\operatorname{Alb} X_\eta$,
alors \uncertain{par} Néron on a

\begin{tikzcd}
X' \arrow[rr, "\alpha"] \arrow[dr, no head] & & A \arrow[dl, no head] \\
& S &
\end{tikzcd}

où \struck{$\alpha$ est défini} $X' = X - x_0$, \uncertain{où} $x_0$ le pt
singulier de \uncertain{$X'$}.
\note{le feuillet s'arrête là ; le reste est vierge.}

\section{Quasi-fonctions de Néron-Weil}
\note{titre souligné de sa main, en tête du feuillet 34.}

\page{34}
$K$ corps muni d'un ens. « propre » $M_K$ de val. abs., qu'on écrit sous
forme logarithmique $v(x)$. $\overline{K}$ une clôture algébrique,
$M_{\overline{K}}$ l'ens. des val. absolues sur $\overline{K}$ prolongeant
celles de $K$. Si $\Gamma$ est le groupe de Galois de $\overline{K}$ sur $K$,
on a
\[
  M_{\overline{K}} / \Gamma \xrightarrow{\ \sim\ } M_K
\]
\note{« algébrique » est biffé puis rétabli par un soulignement pointillé ;
au-dessus, « séparable » est biffé. Un signe $(*)$ cerclé, écrit juste sous
cette formule, semble la désigner : il y est renvoyé plus bas.}

\subsection{Fonctions réelles sur $X(\overline{K}) \times M_{\overline{K}}$, notations}

Soit $X$ un préschéma sur $K$, \struck{\ill{}}
\[
  \varphi : X(\overline{K}) \times M_{\overline{K}} \to \mathbf{R}
\]
une fonction, notée $\varphi(P, v)$. Supposons que $\varphi$
\uncertain{commute avec} $\Gamma$ :
\[
  \varphi(\sigma P, \sigma v) = \varphi(P, v)
\]
alors \struck{pour tout} si \struck{\ill{}}
$P \in X(\overline{K})^{\Gamma} = X(K^{p^{-\infty}})$ (où $K^{p^{-\infty}}$
la clôture parfaite de $K$) \struck{\ill{}} on a
\[
  \varphi(\sigma P, \sigma v) = \varphi(P, v) \quad \text{pour tt } v, \sigma ,
\]
donc en vertu de $(*)$ $\varphi(P, v)$ ne dépend que de $v|K$. Donc $\varphi$
\uncertain{induit}
\[
  X(K^{p^{-\infty}}) \times M_K \to \mathbf{R}
\]
\ill{} fonction $X(K) \times M_K \to \mathbf{R}$.
\note{dans la première formule, un $X$ est surchargé.}

\page{35}
Le même argument montre que pour toute $K$-extension finie $L$ de
$\overline{K}/K$, $\varphi$ induit
\[
  \varphi_L : X(L) \times M_L \to \mathbf{R} ,
\]
ces fonctions $\varphi_L$ ont une propriété de transitivité évidente
\[
  \varphi_L(P, v|L) \simeq \varphi_{L'}(P, v)
  \qquad \left| \begin{array}{l} P \in X(L) \\ v \in M_{L'} \end{array} \right.
\]
et la donnée de $\varphi$ équivaut au système des $\varphi_L$ satisfaisant
la condition de compatibilité précédente. On \struck{dit} désigne par
\[
  \widehat{\Phi}(X)
\]
[ou bien, en cas d'ambiguïté, $\widehat{\Phi}(X/K, M_K)$] l'ens. de ces
$\varphi$. Notons que $\widehat{\Phi}(X)$ est déterminé à isom. unique près,
indépendamment du choix de la clôture alg. $\overline{K}$. On a, pour une
extension \add{finie} $K'$ de $K$, une application naturelle
\[
  \widehat{\Phi}(X/K) \to \widehat{\Phi}(X_{K'}/K')
\]
\struck{et pour $X$ fixé}, qui est \emph{bijective}
\note{« finie » est ajouté au-dessus de la ligne.}

\page{36}
si $K'/K$ est radicielle, et \struck{\ill{}} plus généralement
\struck{telle que} si $K'/K$ est quasi-galoisienne de groupe $\pi$, on a
\[
  \widehat{\Phi}(X/K) \simeq \widehat{\Phi}(X_{K'}/K')^{\pi} .
\]

\subsection{Fonctions et parties $M_K$-bornées}

\struck{Une fonction \ill{} de $X$}

Soit $E$ un ens. muni d'une application $\pi : E \to M_{\overline{K}}$. Une
fonction $\varphi$ sur $E$ est dite $M_K$-\emph{bornée} (rel. à $\pi$) s'il
existe un $M_K$-\emph{diviseur} $c$ \struck{\ill{}} tel que
\[
  |\varphi(\xi)| \leqslant c(\pi(\xi))
\]
i.e. s'il existe \struck{\ill{}} une partie finie $S$ de $M_K$ telle que
\[
  \varphi(\xi) = 0 \quad \text{si } \pi(\xi)|K \notin S
\]
et si de plus $\varphi$ est \emph{bornée}. \struck{\ill{}}

Soit $E$ \add{$= \bigcup_{v \in M_{\overline{K}}} E_v \times \{v\}$} une
partie de $X(\overline{K}) \times M_{\overline{K}}$, on dit que $E$ est
$M_K$-bornée s'il \struck{\ill{}} est contenue dans une réunion
\struck{\ill{}} finie d'ensembles de la forme
\[
  E(U; x_1, \ldots, x_n; c) ,
\]
où $U$ est un ouvert
\marginal{$X$ loc. de t.f. sur $K$}
\note{l'expression de $E$ comme réunion est écrite au-dessus de la ligne ;
la note marginale, en oblique, est rattachée à « Soit ».}

\page{37}
affine de $X$, $\underline{x} = (x_1, \ldots, x_n)$ est un syst. de
générateurs de son anneau affine sur $K$, $c$ un $M_K$-diviseur, et
$E(\ \text{---}\ )$ est l'ens. des
$(P, v) \in X(\overline{K}) \times M_{\overline{K}}$ tels que
\[
  v(\underline{x}(P)) = \operatorname*{Inf}_{1 \leqslant i \leqslant n}
  v(x_i(P)) \geqslant c(v)
\]
[i.e. \add{essentiellement} il existe une partie finie $S$ de $M_K$ telle
que $E$ est formé des $(P, v)$ \add{$\in U(\overline{K}) \times M_{\overline{K}}$}
tels que
\begin{itemize}
  \item[a)] $P$ \uncertain{entier} pour $\underline{x}$, $v$ si
  $v|K \in M_K - S$
  \item[b)] $v(\underline{x}(P)) \geqslant -M$ si $v|K \in S$ ].
\end{itemize}
\note{« essentiellement » et « $\in U(\overline{K}) \times M_{\overline{K}}$ »
sont écrits au-dessus de la ligne.}

\marginal{\emph{N.B.} La notion de partie bornée est inv. par ext. finie du
corps de base. La fonctorialité pour $X \to Y$ \uncertain{n'est pas claire},
aussi il y aurait \uncertain{peut-être} lieu de remplacer $U \subset X$ par
$U \to X$ \ill{}. \emph{Question} : Dans le cas « global », y a-t-il une
notion « $M_K$-bornée » \ill{} de $\mathbf{P}^r_K$, \uncertain{celles} des
\emph{hauteurs} bornées ?}
\note{note marginale, écrite en oblique dans la marge gauche et reliée par
un trait aux conditions a), b).}

\subsection{Fonctions admissibles}

Une fonction $\varphi : X(\overline{K}) \times M_{\overline{K}} \to \mathbf{R}$
est dite \emph{admissible} si elle est \struck{\ill{}} $M_K$-bornée sur toute
partie $M_K$-bornée. Comme les ens.
\[
  \{P\} \times M_{\overline{K}} , \quad \text{où } P \in X(\overline{K}) ,
\]
sont $M_K$-bornés, on voit que \struck{si} pour tt
$P \in X(\overline{K})$ \uncertain{fixé}, la fonction
\[
  \widetilde{\varphi}(P) = \bigl(v \mapsto \varphi(P, v)\bigr) :
  M_{\overline{K}} \to \mathbf{R}
\]
a la \uncertain{vertu} d'être \struck{\ill{}} \uncertain{majorée}
\add{\uncertain{en module}} par une fonction de la forme $c(v)$, avec

\page{38}
$c$ un $M_K$-diviseur positif. D'ailleurs, si $P$ est rationnel sur
\add{l'ext. finie} $K'$ de $K$, alors $\widetilde{\varphi}(P)$ est rat.$/K'$,
i.e. se factorise en
\[
  M_{\overline{K}} \to M_{K'} \to \mathbf{R} ,
\]
où $M_{K'} \to \mathbf{R}$ est induit par un $K'$-diviseur. Donc [posant
\[
  \operatorname{Div}(M_{\overline{K}})
  = \varinjlim_{K'/K} \operatorname{Div}(M_{K'})
\]
$K'$ ext. finie de $\overline{K}/K$] on trouve une application
\[
  \widetilde{\varphi} : X(\overline{K}) \to \operatorname{Div}(M_{\overline{K}})
\]
qui a \ill{} au moins les vertus suivantes
\begin{itemize}
  \item[a)] $\widetilde{\varphi}$ est compatible avec les op. de
  $\Gamma = \operatorname{Gal} \overline{K}/K$
  \item[b)] $\widetilde{\varphi}$ transforme tte partie $M_K$-bornée $F$ de
  $X(\overline{K})$ [i.e. une partie $F$ telle que
  $F \times M_{\overline{K}}$ soit $M_K$-bornée] en une partie \emph{bornée}
  de $M_{\overline{K}}$-diviseurs [bornée au sens de la relation d'ordre, au
  sens : $\exists\, c \in \operatorname{Div}(M_K)$ avec
  $|\widetilde{\varphi}(P)| \leqslant c$ pour tt $P \in F$].
\end{itemize}
\marginal{(Supposons que $\varphi$ est rat.$/K$, i.e. \uncertain{invariante}
par $\operatorname{Gal} \overline{K}/K = \Gamma$)}
\note{« l'ext. finie » est ajouté au-dessus de la ligne. La note marginale
est entre parenthèses en haut à droite du feuillet et se poursuit, cerclée,
en haut à gauche ; elle est rattachée à « D'ailleurs ».}

\page{39}
\subsection{Exemples de $\varphi$ admissibles}

Soit $f \in \Gamma(X, \mathcal{O}_X^*)$, posons
\[
  \varphi_f(P, v) = v(f(P)) \qquad \text{si } P \in X(\overline{K}),\
  v \in M_{\overline{K}} .
\]
C'est une fonction \emph{invariante par $\Gamma$}, je dis qu'elle est
\emph{admissible}. Il suffit de \ill{}

\ill{} que si $X$ est affine, \add{$c \in \operatorname{Div}(M_K)$},
\struck{\ill{}} $K[X]$ engendré par $x_1, \ldots, x_n$, alors $\varphi_f$
est $M_K$-bornée sur $E[X; x_1, \ldots, x_n; c]$, i.e.
$\exists\, c' \in \operatorname{Div}(M_K)$ tel que
\[
  \operatorname*{Inf}_i v(x_i(P)) \geqslant c(v) \Longrightarrow
  |v(f(P))| \leqslant c'(v) .
\]
\note{ce passage, depuis « \ill{} que si $X$ est affine », est encadré et
annulé par deux traits obliques ; la dernière formule est de plus barrée
d'un trait horizontal.}

Par fonctorialité, on est ramené au cas où $X = \mathbf{G}_{m,K}$,
\add{$f$ la section canonique $x$,} \struck{et à vérifier \ill{}}. Dans ce
cas, \struck{pour tt} une partie $M_K$-bornée de
$X(\overline{K}) \times M_{\overline{K}}$
\add{$= \overline{K}^* \times M_{\overline{K}}$} est contenue dans une partie
de la forme $E(M, S)$, $S$ partie finie de $M_K$, \struck{$M \geqslant 0$},
$M$ un nb $> 0$, formée des $(x, v)$ tels que
\begin{itemize}
  \item[a)] $v(x) = 0$ \struck{\ill{}} \add{i.e. $x$ unité} si
  $v|K \notin S$
  \item[b)] \struck{\ill{}} $|v(x)| \leqslant M$ si $v|K \in S$.
\end{itemize}
\note{« $f$ la section canonique $x$ », « $= \overline{K}^* \times M_{\overline{K}}$ » et « i.e. $x$ unité » sont écrits au-dessus des lignes.}

\page{40}
Or on a
\[
  \varphi_f(x, v) = v(x)
\]
\struck{\ill{}} d'où aussitôt la conclusion.

\subsection{Faisceautisation}

Pour tt ouvert $U$ de $X$, soit $\underline{\Phi}_X(U)$ l'ens. des fonctions
admissibles $\Gamma$-invariantes sur $U(\overline{K}) \times M_{\overline{K}}$.
Pour $U$ variable, on trouve ainsi un \emph{préfaisceau}
$\underline{\Phi}_X$ sur $X$, je dis que c'est un \emph{faisceau}.
Évidemment, c'est un préfaisceau « séparé », reste à montrer que si une
fonction
\[
  \varphi : X(\overline{K}) \times M_{\overline{K}} \to \mathbf{R}
\]
invariante par $\Gamma$ est telle que l'on ait
$\varphi | U_i(\overline{K}) \times M_{\overline{K}}$ admissible pour tt
$i$ (où $(U_i)$ est un recouvrement ouvert de $X$), alors $\varphi$ est
admissible. On est ramené au cas où $X$ est affine, et il semble qu'on ait
besoin de savoir que \struck{tte} \add{\uncertain{la réunion}}
\emph{partie} $M_K$-\emph{bornée de} $X(\overline{K}) \times M_{\overline{K}}$
\note{le feuillet s'arrête au milieu de la phrase ; « la réunion » est écrit
au-dessus de « tte », biffé. La lettre de $\underline{\Phi}_X$ est celle de
$\widehat{\Phi}$ (feuillets 35-36), soulignée ; lue $\Phi$ sous réserve.}

\end{document}
